\documentclass[11pt]{article}
\usepackage[margin=0.9in]{geometry}
\usepackage{amsmath,amssymb,amsthm,mathtools}
\usepackage[T1]{fontenc}
\usepackage[utf8]{inputenc}
\usepackage{enumitem}
\usepackage{booktabs,longtable,array,seqsplit}
\usepackage{hyperref}
\usepackage{xcolor}
\usepackage{microtype}
\hypersetup{colorlinks=true,linkcolor=blue!45!black,citecolor=blue!45!black,urlcolor=blue!45!black}
\setlength{\emergencystretch}{6em}

\newtheorem{theorem}{Theorem}[section]
\newtheorem{proposition}[theorem]{Proposition}
\newtheorem{lemma}[theorem]{Lemma}
\newtheorem{corollary}[theorem]{Corollary}
\newtheorem{claim}[theorem]{Claim}
\theoremstyle{definition}
\newtheorem{definition}[theorem]{Definition}
\newtheorem{assumption}[theorem]{Assumption}
\theoremstyle{remark}
\newtheorem{remark}[theorem]{Remark}

\newcommand{\R}{\mathbb{R}}
\newcommand{\C}{\mathbb{C}}
\newcommand{\N}{\mathbb{N}}
\newcommand{\Hsp}{\mathcal{H}}
\newcommand{\Fsp}{\mathcal{F}}
\newcommand{\Id}{\mathrm{Id}}
\newcommand{\tr}{\operatorname{tr}}
\newcommand{\diag}{\operatorname{diag}}
\newcommand{\rank}{\operatorname{rank}}
\newcommand{\spec}{\operatorname{spec}}
\newcommand{\spanop}{\operatorname{span}}
\newcommand{\dist}{\operatorname{dist}}
\newcommand{\range}{\operatorname{ran}}
\newcommand{\Ker}{\operatorname{ker}}
\newcommand{\Pinch}{\mathcal{C}}
\newcommand{\Off}{\widetilde{\mathcal{C}}}
\newcommand{\Pperp}{P^{\perp}}
\newcommand{\Qperp}{Q^{\perp}}
\newcommand{\op}{\mathrm{op}}
\newcommand{\F}{\mathrm{F}}
\newcommand{\HS}{\mathrm{HS}}
\newcommand{\ip}[2]{\left\langle #1,#2\right\rangle}
\newcommand{\norm}[1]{\left\lVert #1\right\rVert}
\newcommand{\abs}[1]{\left\lvert #1\right\rvert}
\newcommand{\code}[1]{\nolinkurl{#1}}

\newenvironment{argument}{\paragraph{Argument.}\begin{enumerate}[leftmargin=2em]}{\end{enumerate}}
\newenvironment{proofoutline}{\paragraph{Proof architecture.}\begin{enumerate}[leftmargin=2em]}{\end{enumerate}}
\newcommand{\sourceanchor}[1]{\par\smallskip\noindent\textbf{Source anchor.} #1\par\smallskip}
\newcommand{\sourcecomment}[1]{\begin{remark}#1\end{remark}}
\newenvironment{sourceguide}{\begin{description}[style=nextline,leftmargin=3.3cm,labelwidth=3.1cm]}{\end{description}}
\newenvironment{formalizationmap}{\begin{longtable}{@{}p{0.20\textwidth}p{0.31\textwidth}p{0.40\textwidth}@{}}\toprule Source item & Modern mathematical content & Repository status \\\midrule\endfirsthead\toprule Source item & Modern mathematical content & Repository status \\\midrule\endhead}{\bottomrule\end{longtable}}

\title{Kostrykin--Makarov--Motovilov (2005): Riccati graph subspaces and the tan $\Theta$ theorem}
\author{}
\date{Distilled reconstruction for the finite Davis--Kahan formalization}
\begin{document}
\maketitle

\section*{Source map and formalization scope}
\begin{sourceguide}
\item[Primary source] Vadim Kostrykin, Konstantin A. Makarov, and Alexander K. Motovilov, \emph{On the Existence of Solutions to the Operator Riccati Equation and the tan $\Theta$ Theorem}, Integral Equations and Operator Theory 51 (2005), 121--140; preprint \href{https://arxiv.org/abs/math/0210032}{arXiv:math/0210032}.
\item[Coverage] Theorems 1 and 2; graph-subspace equivalence and block diagonalization in Theorem 2.3; spectral uniqueness; the a priori estimate of Theorem 5.1; the use of tan $2\Theta$ for strict contractivity.
\item[Formalization target] \code{DavisKahan.Experimental.Literature.Riccati}, \code{DavisKahanExt.Riccati}, and the graph/no-pole interpretation of the source-faithful finite tan $\Theta$ theorem.
\item[Scope boundary] The paper permits a possibly unbounded diagonal block $C$ and uses analytic factorization for existence. The repository's immediate target is finite-dimensional bounded operators over arbitrary \code{RCLike}; the graph, Riccati, angle, uniqueness, and norm arguments survive without domain theory.
\item[Logical separation] The paper has three distinct layers: existence of a graph solution, identification of its spectral location, and a tan bound once those facts hold. Do not collapse them into one theorem with hidden assumptions.
\end{sourceguide}

\section{The block operator and its Riccati equation}

On an orthogonal decomposition
\[
 \Hsp=\Hsp_A\oplus\Hsp_C,
\]
consider
\[
 H=
 \begin{pmatrix}
 A&B\\
 B^*&C
 \end{pmatrix},
\]
where $A=A^*$ is bounded, $C=C^*$ may be unbounded in the source, and $B:\Hsp_C\to\Hsp_A$ is bounded. A bounded operator $X:\Hsp_A\to\Hsp_C$ is sought satisfying
\begin{equation}
 XA-CX+XBX=B^*.
 \tag{KMM-R}\label{KMM-R}
\end{equation}
The graph is
\[
 \mathcal G(X)=\{x\oplus Xx:x\in\Hsp_A\}.
\]

The nonlinear term $XBX$ is not an optional complication. It records the fact that the target subspace is invariant for the full block operator rather than merely approximately invariant for its diagonal part.

\section{Graph reduction if and only if Riccati}

For $x\in\Hsp_A$,
\[
 H(x\oplus Xx)
 =
 (Ax+BXx)\oplus(B^*x+CXx).
\]
This belongs to $\mathcal G(X)$ exactly when
\[
 B^*x+CXx=X(Ax+BXx),
\]
which is equivalent to \eqref{KMM-R}. Since $H$ is self-adjoint, invariance of the graph implies reduction.

The paper packages the equivalence as:
\[
 X\text{ solves \eqref{KMM-R}}
 \quad\Longleftrightarrow\quad
 \mathcal G(X)\text{ reduces }H.
\]
This elementary finite-dimensional identity should be a foundational reusable theorem, not reproved inside every tan argument.

Define
\[
 V=
 \begin{pmatrix}
 I&-X^*\\
 X&I
 \end{pmatrix}.
\]
When $X$ solves the Riccati equation,
\[
 V^{-1}HV
 =
 \begin{pmatrix}
 A+BX&0\\
 0&C-B^*X^*
 \end{pmatrix}.
\]
The diagonal blocks are generally only similar to self-adjoint operators. Conjugating them by $(I+X^*X)^{1/2}$ and $(I+XX^*)^{1/2}$ produces self-adjoint representatives. This explains why their spectra are real and why they are the correct spectral parts attached to the graph and its orthogonal complement.

\section{The graph-angle dictionary}

Let $P_A$ be the projection onto $\Hsp_A\oplus\{0\}$ and $P_X$ the projection onto $\mathcal G(X)$. The operator angle $\Theta$ between these subspaces satisfies
\[
 \Theta=\arctan\sqrt{X^*X}.
\]
Consequently
\[
 \tan\Theta=\sqrt{X^*X},
 \qquad
 \norm{\tan\Theta}_{\op}=\norm{X}_{\op},
\]
and
\[
 \norm{P_A-P_X}_{\op}
 =\frac{\norm X}{\sqrt{1+\norm X^2}}.
\]
Thus strict contractivity $\norm X<1$ is exactly the acute branch condition
\[
 \norm\Theta<\frac{\pi}{4}.
\]

The finite formalization should expose this dictionary in both directions:
\begin{itemize}[leftmargin=2em]
\item transversality to the vertical coordinate space gives a unique graph operator;
\item the graph operator gives the tangent-angle operator;
\item contractivity selects the quarter-turn branch required by tan $2\Theta$.
\end{itemize}

\section{Theorem 1: existence in a finite gap}

Assume $C$ has a finite spectral gap
\[
 \Delta=(\alpha,\beta)
\]
containing $\spec(A)$, and let
\[
 d=\dist(\spec(A),\spec(C)).
\]
Theorem 1(i) states that if
\[
 \norm B<\sqrt{d\abs\Delta},
\]
then the spectrum of $H$ inside $\Delta$ is a proper closed subset of $\Delta$, and its spectral subspace is the graph of a bounded solution $X$ of \eqref{KMM-R}. Moreover $X$ is uniquely characterized among bounded solutions by
\[
 \spec(A+BX)\subset\Delta,
 \qquad
 \spec(C-B^*X^*)\subset\R\setminus\Delta.
\]

If $C$ is also bounded and
\[
 \norm B<\sqrt{d(\abs\Delta-d)},
\]
then Theorem 1(ii) proves $\norm X<1$ and gives the half-angle estimate
\[
 \norm X
 \le
 \tan\left[
  \frac12\arctan
  \frac{2\norm{AB+BC}}
       {d(\abs\Delta-d)-\norm B^2}
 \right]
 <1.
\]
The exact expression belongs to the paper's a priori existence theory. It is not the same theorem as the simpler a posteriori tan bound below.

\section{Theorem 2: the extended tan $\Theta$ estimate}

Theorem 2 makes no a priori assumption on the mutual position of $\spec(A)$ and $\spec(C)$. Instead it assumes that a bounded Riccati solution already exists and that its graph spectral part lies inside a gap of $C$.

Let $\Delta$ be a finite or infinite spectral gap of $C$. Suppose $X$ solves \eqref{KMM-R}, so $\mathcal G(X)$ reduces $H$, and suppose
\[
 \spec(H|_{\mathcal G(X)})
\]
is a proper closed subset of $\Delta$. Set
\[
 \delta
 =\dist\bigl(\spec(H|_{\mathcal G(X)}),\spec(C)\bigr).
\]
Then
\begin{equation}
 \norm X\le\frac{\norm B}{\delta},
 \qquad
 \norm{\tan\Theta}\le\frac{\norm B}{\delta}.
 \tag{KMM-tan}\label{KMM-tan}
\end{equation}

This is the decisive source-faithful denominator. It uses the spectrum of the actual graph restriction $H|_{\mathcal G(X)}$, equivalently the self-adjoint representative of $A+BX$, against the original complementary block $C$.

The paper explicitly warns that a tempting symmetric replacement generally fails. One cannot simply replace $\delta$ by
\[
 \dist\bigl(\spec(C-B^*X^*),\spec(A)\bigr)
\]
and expect the same inequality. The direction of the spectral gap is part of the theorem.

\section{Why the pole cannot occur}

The graph formulation itself excludes an angle of $\pi/2$: $\mathcal G(X)$ is transverse to $\Hsp_C$. In a theorem stated directly for two subspaces, this must be proved rather than assumed.

A finite-dimensional route is:
\begin{enumerate}[leftmargin=2em]
\item identify the target reducing subspace by its spectral location;
\item prove its intersection with the wrong coordinate subspace is zero using the spectral gap;
\item use dimension balance to obtain transversality in both directions;
\item construct the unique graph map $X$;
\item translate its norm to $\tan\Theta$.
\end{enumerate}
This is the bridge between the abstract Riccati theorem and Nakatsukasa's pole-free finite matrix statement.

\section{Proof mechanism for the tan bound}

After conjugating the graph restriction to a self-adjoint operator $\Lambda$ on $\Hsp_A$, the Riccati equation can be reorganized as a linear Sylvester equation for a normalized graph operator. Schematically,
\[
 C Y-Y\Lambda=T,
\]
where the normalization absorbs $(I+X^*X)^{\pm1/2}$ and $T$ is controlled by $B$.

The spectra of $C$ and $\Lambda$ are separated by $\delta$. The ordered/interval gap permits a constant-one Sylvester estimate,
\[
 \delta\norm Y\le\norm T.
\]
Undoing the normalization yields \eqref{KMM-tan}. This is why the result remains sharp and why the proof does not require a smallness assumption on $B$ once the graph and spectral localization are given.

For Lean, there are two possible proof presentations:
\begin{itemize}[leftmargin=2em]
\item derive the exact normalized Sylvester equation and apply the existing ordered-gap theorem; or
\item use the current coercive per-vector tan argument directly, then prove equivalence to the graph formulation.
\end{itemize}
The first is closer to the paper's operator architecture; the second may be shorter for the existing finite API.

\section{A priori bounds and strict contractivity}

Theorem 5.1 combines existence and spectral localization with Theorem 2. If the graph solution from Theorem 1 has its perturbed spectral set at distance $\widetilde\delta$ from $\spec(C)$, then
\[
 \norm X\le\frac{\norm B}{\widetilde\delta}.
\]
The work is therefore split cleanly:
\begin{enumerate}[leftmargin=2em]
\item construct the graph spectral subspace;
\item localize its spectrum;
\item feed that distance into the a posteriori tan theorem.
\end{enumerate}

For strict contractivity, the paper shifts the block operator and applies a tan $2\Theta$ theorem to its square. Squaring converts a finite interior gap into an ordered spectral configuration while preserving an off-diagonal structure. The resulting half-angle formula proves $\norm X<1$ under the stated threshold. A low-dimensional example shows the threshold is sharp.

This shifted-square idea is valuable for the finite formalization: branch selection for tan $2\Theta$ can often be proved from a positivity or ordered-gap theorem applied to a square, instead of manipulating inverse cosine operators directly.

\section{Finite formalization map}

\begin{formalizationmap}
Block operator data & $H=\begin{psmallmatrix}A&B\\B^*&C\end{psmallmatrix}$ on an orthogonal sum & \code{BlockOperatorData} and \code{blockOperator} are scaffolded.\\
Graph subspace & $\mathcal G(X)=\{x\oplus Xx\}$ & \code{blockGraph} exists.\\
Riccati equation & Graph invariance/reduction iff nonlinear defect vanishes & \code{riccatiDefect}, \code{SolvesRiccati}, and \code{graph_reduces_iff_solvesRiccati} exist.\\
Block diagonalization & Similarity to $A+BX$ and $C-B^*X^*$ & A theorem seam exists but should be connected to spectral localization.\\
Graph angle & $\Theta=\arctan\sqrt{X^*X}$ and $\norm{\tan\Theta}=\norm X$ & Needs a stable coordinate-free theorem connecting graph maps to existing angle operators.\\
A posteriori tan theorem & Gap between graph restriction and $C$ gives $\delta\norm X\le\norm B$ & The finite per-vector tan theorem captures the no-pole content; a graph wrapper remains high value.\\
Existence and uniqueness & Finite-gap spectral subspace is the unique graph solution with specified spectral parts & Scaffolded but not the immediate prerequisite for the classical Part III endpoint.\\
Strict contraction & Shift-square plus tan $2\Theta$ selects $\norm X<1$ & Maps to quarter-turn avoidance and canonical branch selection.\\
\end{formalizationmap}

A high-value implementation order is:
\begin{enumerate}[leftmargin=2em]
\item finish the graph/reduction/Riccati equivalence in the reusable bounded finite setting;
\item prove the graph-angle and projector formulas;
\item wrap the existing no-pole tan theorem as a graph-subspace theorem with the exact KMM denominator;
\item prove spectral-subspace transversality automatically under the finite interval/exterior hypotheses;
\item only then attack global existence or analytic factorization analogues.
\end{enumerate}

\section{Failure modes to exclude}

\begin{itemize}[leftmargin=2em]
\item Do not infer existence of a Riccati solution from spectral separation alone; the paper stresses that separation by itself is insufficient in general.
\item Do not use the spectrum of $C-B^*X^*$ against $A$ as an interchangeable denominator.
\item Do not state a tan theorem before proving transversality or otherwise excluding a principal angle of $\pi/2$.
\item Do not confuse the a posteriori estimate \eqref{KMM-tan} with the small-perturbation assumptions needed to construct the graph in Theorem 1.
\item Do not treat strict contractivity as a mere numerical improvement; it is the branch-selection statement $\Theta<\pi/4$.
\end{itemize}

\section*{Distilled conclusion}

The Riccati equation is the algebraic certificate that a perturbed spectral subspace is a graph. Once that graph and its spectral location are known, the tan $\Theta$ estimate is a sharp constant-one Sylvester theorem with no perturbation-smallness assumption. The finite Lean development should separate existence, transversality, graph-angle geometry, and the final norm bound, then expose thin wrappers connecting the source theorem to the classical Part III API.

\end{document}
